\section{Additional Comparative Statics and Counterexamples}
\label{supp:additional-comparative-statics}\label{app:counterexamples}

This section records the supporting one-way comparative statics and the exact
counterexamples that delimit the coverage results in
\cref{supp:belief-occupation-coverage}.  Every comparison holds conditional
on the displayed current state and fixed action or on the stated one-shot
problem.  Unnamed primitives are held fixed.  These results do not promote a
one-shot cost-covering region to an optimized Bellman research region.

\subsection{Cost, admission, survival, and delay}

The supporting order comparisons are compact because every nondisplayed
primitive is held fixed:
\begin{center}
\small
\renewcommand{\arraystretch}{1.08}
\begin{tabular}{@{}>{\raggedright\arraybackslash}p{0.18\textwidth}
                    >{\raggedright\arraybackslash}p{0.25\textwidth}
                    >{\raggedright\arraybackslash}p{0.43\textwidth}@{}}
\toprule
Change & Direction & Required boundary \\
\midrule
higher initiation cost
  & optimized value weakly falls
  & the same project actions change; Continue and other primitives are fixed \\
higher admission or survival
  & project value weakly rises
  & successful completion weakly dominates failure \\
longer duration
  & project value weakly falls
  & every added operating date satisfies
    \(F_t\le(1-\beta)W\); suspended operation is \(F_t=0\) \\
higher incumbent frontier
  & fixed-candidate operating gain weakly falls
  & candidate profile and research opportunities are fixed \\
\bottomrule
\end{tabular}
\end{center}

\begin{proposition}[One-shot cutoff comparative statics]
\label{prop:one-shot-cutoff-comparative-statics}
Holding all other one-shot primitives fixed, higher cost shrinks
\(\mathcal C\); higher nonnegative survival or admission expands it; and a
higher operating frontier shrinks it for a fixed candidate.  When both
compared nonempty regions have upper-threshold representations, higher cost
or frontier raises the cutoff, while higher survival or admission lowers it.
\end{proposition}

These set inclusions concern the one-shot cost-covering region in
\cref{def:single-gap-region}.  Cost and survival also order an exact finite
weak-research region for a fixed research action versus Continue when all
other unified-model returns are unchanged.  Neither result, by itself, orders
the full optimal research region when opportunity menus, continuation laws,
or competing actions also change.

The cutoff directions are pointwise set comparisons.  Raising \(\kappa\)
removes weakly covered beliefs.  Raising a nonnegative survival or admission
factor raises \(G\) and adds beliefs.  Raising the frontier lowers the
positive-part gap and removes beliefs.  If two nonempty regions are
\(\operatorname{Ici}(b_0^*)\) and \(\operatorname{Ici}(b_1^*)\), set
inclusion reverses the cutoff order, yielding the directions in
\cref{prop:one-shot-cutoff-comparative-statics}.

At a fixed finite horizon and conditional on the current decision state,
holding every nondisplayed primitive fixed, a higher initiation cost subtracts
the same nonnegative increment from each affected project action and leaves
Continue unchanged.  Backward induction and maximization over the finite
action menu therefore make optimized value antitone in the cost schedule.

For a binary completion law, condition on the fixed current state and chosen
project, with success continuation \(W^+\) and failure continuation \(W^-\).
The binary specialization assumes the conditional success mass
\(m=\pi\rho^d\); it is not inferred from unrelated marginals.  The completion
component is
\[
  mW^+ +(1-m)W^-.
\]
For two admission--survival configurations its exact difference is
\[
  \bigl[m_1W^+ +(1-m_1)W^-\bigr]
  -
  \bigl[m_0W^+ +(1-m_0)W^-\bigr]
  =(m_1-m_0)(W^+-W^-).
\]
Thus greater admission or survival raises project value when
\(W^+\ge W^-\), and reverses sign if success is harmful.

For continued operation, keep the initial state, belief path, and terminal
continuation \(W\) fixed while comparing durations.  Write the elapsed return
at duration \(d\) as
\[
  R_d=-\kappa+\sum_{t<d}\beta^tF_t+\beta^dW .
\]
The one-date increment is
\[
  R_{d+1}-R_d
  =\beta^d\!\left[F_d-(1-\beta)W\right].
\]
Consequently, for \(d_0<d_1\),
\[
  R_{d_1}-R_{d_0}
  =\sum_{t=d_0}^{d_1-1}
    \beta^t\!\left[F_t-(1-\beta)W\right].
\]
Longer delay is therefore harmful when every added date satisfies
\(F_t\le(1-\beta)W\).  Suspended operation is the special case \(F_t=0\).
Nonnegativity alone is insufficient: with \(\beta=1/2\) and \(F_t=W=1\),
\(R_1=3/2<R_2=7/4\).

\subsection{Threshold and multi-gap counterexamples}

The assumptions of the positive threshold theorem have exact finite
boundaries.  The following examples preserve the original rational values
and isolate the failed condition in each case.

\begin{counterexample}[A nonmonotone kernel defeats the upper threshold]
\label{cx:single-peaked-disconnected}
On \(B=\{0,1,2\}\), take deterministic destinations \((1,0,1)\).
The resulting row-stochastic kernel maps the increasing nonnegative gap
\(\Delta=(0,1,2)\) to \(P\Delta=(1,0,1)\).  With unit survival and
antitone constant unit cost, the weak cost-covering set is \(\{0,2\}\), which
is neither upper nor order-connected.  Thus the stochastic-monotonicity
assumption is independently necessary.  The same kernel also maps the
single-peaked gap \((0,1,0)\) to \((1,0,1)\), retaining the earlier boundary
against unrestricted single-peaked preservation.
\end{counterexample}

\begin{counterexample}[Non-antitone cost disconnects a monotone covering set]
\label{cx:arbitrary-cost-disconnected}
The increasing potential \((1,2,3)\) and cost \((0,3,0)\) produce the
cost-covering set \(\{1,3\}\).  Antitonicity of research cost cannot be dropped
from \cref{thm:finite-monotone-coverage}.
\end{counterexample}

\begin{counterexample}[One-project multi-gap disconnection]
\label{cx:multi-gap-disconnection}
Let \(B=\{0,1,2,3,4\}\).  One project supplies candidates with certified gaps
\[
  \Delta^L=(4,0,0,0,0),\qquad
  \Delta^R=(0,0,0,0,4)
\]
over the zero frontier.  Under the degree-four Bernstein/binomial kernel
\[
  P_{ij}=\binom{4}{j}(i/4)^j(1-i/4)^{4-j},
  \qquad i,j\in B,
\]
their aggregate one-shot coverage potential is
\(P(\Delta^L+\Delta^R)=(4,41/32,1/2,41/32,4)\).  Constant strict cost one
therefore gives the one-shot cost-covering set
\(\{0,1\}\cup\{3,4\}\), which is not order-connected.
\end{counterexample}

\begin{counterexample}[Arbitrary cost defeats a component bound]
\label{cx:arbitrary-cost-components}
With the same five-state kernel, constant gap and potential
\((2,2,2,2,2)\), and cost \((0,3,0,3,0)\), the strict cost-covering set is
\(\{0,2,4\}\).  A kernel-only variation-diminishing condition cannot control
net-region components under unrestricted belief-dependent cost.
\end{counterexample}

The multi-gap example shows that the upper-threshold conclusion does not
extend to separated candidate gaps.  Separately, the first two examples show
that neither monotone gaps alone nor a well-behaved gross-value profile alone
ensures connected cost covering.  The arbitrary-cost example rules out a
kernel-only component bound when belief-dependent cost is unrestricted.
